\documentclass[11pt,a4paper]{article}

% ---------------------------
% Core Packages
% ---------------------------
\usepackage[top=1in, bottom=1in, left=1in, right=1in]{geometry}
\usepackage{amsmath,amssymb,amsthm}
\usepackage{mathptmx}
\usepackage{siunitx}
\usepackage{graphicx}
\usepackage{booktabs}
\usepackage{authblk}
\usepackage[colorlinks=true,linkcolor=black,citecolor=blue,urlcolor=blue]{hyperref}
\usepackage{url}
\usepackage{enumitem}
\usepackage{caption}
\usepackage{subcaption}
\usepackage{float}
\usepackage{listings}
\usepackage{xcolor}
\usepackage{titlesec}
\usepackage{fancyhdr}
\usepackage{tcolorbox}

% ---------------------------
% Watermark Configuration
% ---------------------------
\usepackage{draftwatermark}
\SetWatermarkText{UNPUBLISHED DRAFT - CONFIDENTIAL}
\SetWatermarkScale{0.7}
\SetWatermarkAngle{55}
\SetWatermarkColor[gray]{0.85}
\SetWatermarkFontSize{2cm}

% ---------------------------
% Spacing Configuration
% ---------------------------
\setlength{\parskip}{6pt}
\setlength{\parindent}{0pt}
\renewcommand{\baselinestretch}{1.08}

% Section spacing
\titlespacing*{\section}{0pt}{14pt plus 2pt minus 2pt}{8pt plus 1pt minus 1pt}
\titlespacing*{\subsection}{0pt}{12pt plus 2pt minus 1pt}{6pt plus 1pt minus 1pt}
\titlespacing*{\subsubsection}{0pt}{10pt plus 1pt minus 1pt}{4pt plus 1pt minus 1pt}

% Section formatting
\titleformat{\section}
  {\normalfont\fontsize{13}{15}\bfseries\sffamily}{\thesection}{1em}{}
\titleformat{\subsection}
  {\normalfont\fontsize{11.5}{14}\bfseries\sffamily}{\thesubsection}{1em}{}
\titleformat{\subsubsection}
  {\normalfont\fontsize{11}{13}\bfseries\itshape}{\thesubsubsection}{1em}{}

% ---------------------------
% Caption Configuration
% ---------------------------
\captionsetup{
  font={small,sf},
  labelfont={bf,sf},
  format=plain,
  justification=justified,
  singlelinecheck=false,
  skip=8pt,
  belowskip=0pt
}

\graphicspath{{figures/}}

% ---------------------------
% Code Listings
% ---------------------------
\definecolor{codebg}{RGB}{248,248,248}
\definecolor{codeblue}{RGB}{0,102,204}
\definecolor{codecomment}{RGB}{102,102,102}
\definecolor{codestring}{RGB}{221,102,0}

\lstdefinestyle{code}{
  basicstyle=\ttfamily\footnotesize,
  keywordstyle=\color{codeblue}\bfseries,
  commentstyle=\color{codecomment}\itshape,
  stringstyle=\color{codestring},
  showstringspaces=false,
  breaklines=true,
  frame=single,
  rulecolor=\color{black!20},
  backgroundcolor=\color{codebg},
  xleftmargin=8pt,
  xrightmargin=8pt,
  framexleftmargin=8pt,
  framexrightmargin=8pt,
  numbers=left,
  numberstyle=\tiny\color{black!40},
  numbersep=10pt,
  aboveskip=10pt,
  belowskip=10pt
}
\lstset{style=code}

% ---------------------------
% SI Units
% ---------------------------
\sisetup{
  detect-all,
  group-separator={,},
  detect-weight=true,
  detect-family=true
}

% ---------------------------
% Header/Footer
% ---------------------------
\pagestyle{fancy}
\fancyhf{}
\fancyhead[L]{\small\sffamily\textit{LunarAtlas: Chandrayaan-3 LIBS Infrastructure}}
\fancyhead[R]{\small\sffamily\textit{Anand \& Saeed, 2026}}
\fancyfoot[C]{\small\sffamily\thepage}
\renewcommand{\headrulewidth}{0.3pt}
\renewcommand{\footrulewidth}{0pt}

\fancypagestyle{firstpage}{
  \fancyhf{}
  \renewcommand{\headrulewidth}{0pt}
  \fancyfoot[C]{\small\sffamily\thepage}
}

% ---------------------------
% Custom environments
% ---------------------------
\setlist[enumerate]{topsep=4pt,itemsep=3pt,partopsep=0pt,parsep=0pt,leftmargin=*}
\setlist[itemize]{topsep=4pt,itemsep=3pt,partopsep=0pt,parsep=0pt,leftmargin=*}
\setlist[description]{topsep=4pt,itemsep=3pt,partopsep=0pt,parsep=0pt,leftmargin=0pt}

% ---------------------------
% Title Setup
% ---------------------------
\title{\vspace{-8mm}\sffamily\bfseries\LARGE
LunarAtlas: A Reproducible Spectral Data Infrastructure for Chandrayaan-3 LIBS with Peak-Preserving Visualization
\vspace{2mm}}

\author[1]{Lovekesh Anand}
\author[1]{Dua Saeed}

\affil[1]{\small\sffamily Independent Researchers, India}
\affil[ ]{\vspace{2mm}}

\date{\sffamily April 2026}

% ---------------------------
% Document
% ---------------------------
\begin{document}

\maketitle
\thispagestyle{firstpage}
\vspace{-6mm}

% ---------------------------
% Abstract Box
% ---------------------------
\begin{tcolorbox}[
  colback=black!3,
  colframe=black!30,
  boxrule=0.5pt,
  arc=2pt,
  left=8pt,
  right=8pt,
  top=8pt,
  bottom=8pt,
  before skip=8pt,
  after skip=12pt
]
\small\sffamily
\textbf{Abstract.} We present \textbf{LunarAtlas}, a backend-centric scientific data infrastructure that transforms publicly available Chandrayaan-3 Laser-Induced Breakdown Spectroscopy (LIBS) Level-1 data into a reproducible, versioned, and machine-accessible spectral database optimized for interactive visualization and scientific analysis. Unlike conventional planetary mission archives that expose static files or generic REST endpoints, LunarAtlas implements all scientifically critical operations—wavelength calibration, spectral cleaning, peak detection, and NIST validation—within a PostgreSQL-backed server architecture, ensuring reproducibility and auditability.

We introduce an \textit{adaptive min-max downsampling algorithm} specifically designed for spectroscopic data that guarantees preservation of emission line peaks while reducing data volume by up to two orders of magnitude for efficient web-based visualization. The algorithm combines zoom-aware bucket sizing, overlapping bucket boundaries to prevent peak splitting, and NIST-referenced peak insertion to maintain scientific fidelity across all zoom levels.

Using processed Chandrayaan-3 LIBS spectra, we demonstrate successful identification of major lunar elements (Mg, Ca, Al, Fe, Si, O) with wavelength accuracies within the instrumental resolution, and sub-\SI{500}{ms} API response times for typical $\sim 1.6\times10^{4}$-point observations. The system is validated on real Level-1 files retrieved from ISRO's public PDS4-compliant archive, establishing a methodological framework for planetary spectroscopy data systems with applications to Mars LIBS missions (ChemCam, SuperCam, MarSCoDe) and future lunar programs.

\vspace{4pt}
\textbf{Keywords:} Laser-Induced Breakdown Spectroscopy • Chandrayaan-3 • Planetary Data Systems • PDS4 • Spectral Visualization • Data Infrastructure • Peak Detection • NIST Validation • Lunar Composition
\end{tcolorbox}

\vspace{2mm}

% ---------------------------
% Main Content
% ---------------------------

\section{Introduction}

\subsection{Background and Motivation}

Laser-Induced Breakdown Spectroscopy (LIBS) has emerged as a transformative technology for \textit{in situ} planetary exploration, enabling rapid, stand-off elemental analysis without sample preparation~\cite{Cremers2006}. The technique operates by focusing a pulsed laser onto a target surface, creating a transient plasma plume whose characteristic emission spectrum encodes the elemental composition of the ablated material. Since its first planetary deployment on NASA's Mars Science Laboratory (MSL) rover \textit{Curiosity} in 2012, LIBS has revolutionized our understanding of Martian geochemistry through the ChemCam instrument~\cite{Wiens2012,Maurice2012}.

ChemCam combines a mast-mounted laser and telescope with body-mounted spectrometers spanning 240--905~nm and has acquired more than 188,000 LIBS spectra during the nominal mission in Gale Crater~\cite{Mangold2016}. SuperCam on the Mars 2020 \textit{Perseverance} rover extends this approach by integrating LIBS with Raman, time-resolved luminescence, and VIS-IR reflectance spectroscopy~\cite{Wiens2020}. China's Tianwen-1 mission carries the Mars Surface Composition Detector (MarSCoDe), combining LIBS and short-wave infrared spectroscopy~\cite{Xu2021}.

The extension of LIBS to lunar exploration represents a significant milestone. India's Chandrayaan-3 mission, which successfully landed near the lunar south pole in August 2023, carried a LIBS instrument aboard the Pragyan rover alongside an Alpha Particle X-ray Spectrometer (APXS)~\cite{ISROCh3}. Operating in the near-vacuum lunar environment presents unique challenges compared to Martian LIBS due to different plasma expansion dynamics, lack of atmospheric confinement, and extreme temperature variations.

\subsection{The Challenge of Planetary Spectroscopy Data Accessibility}

While the scientific value of planetary LIBS data is well established, significant barriers exist between raw mission data and actionable scientific insights. Planetary Data System version~4 (PDS4) archives provide robust standards for long-term data preservation through self-describing XML labels paired with tabular or binary data products~\cite{Hughes2018}. However, PDS4 archives are optimized for archival completeness rather than interactive analysis. Users must manually interpret XML schemas, decode instrument-specific metadata, and reconcile calibration context before performing spectral analysis.

High-resolution LIBS spectra typically contain $\sim 8{,}000$--$20{,}000$ wavelength-intensity pairs per observation, and direct rendering can saturate browser memory and produce sluggish interactions. Published analyses often rely on ad hoc scripts that are not version-controlled or archived, hampering reproducibility. Generic downsampling algorithms (mean/median aggregation, Douglas-Peucker, Largest Triangle Three Buckets) prioritize visual shape over peak preservation and risk loss of narrow emission lines that carry elemental information.

These challenges are particularly acute for Chandrayaan-3 LIBS data which, while publicly released through ISRO's PDS4-compliant archives, lack mission-provided, peak-aware interactive analysis tools comparable to those developed for Mars LIBS missions.

\subsection{Research Objectives and Contributions}

\textit{Research Question:} How can publicly available lunar LIBS data be transformed into a reproducible, versioned, and machine-accessible system without compromising scientific validity or peak fidelity?

This paper makes the following contributions:

\begin{enumerate}[label=\textbf{C\arabic*.},leftmargin=1.5em]
\setlength{\itemsep}{2pt}
\item \textbf{PDS4-aware data pipeline:} Systematic ingestion framework parsing Chandrayaan-3 LIBS XML labels and Level-1 CSV products, extracting wavelength-intensity spectra into normalized PostgreSQL schema preserving scientific hierarchy.

\item \textbf{Physically motivated cleaning:} Background-subtraction and noise-rejection methodology tailored to lunar LIBS, converting wide L1 tables into long-form records while preserving spectral line positions and relative intensities.

\item \textbf{Adaptive min-max downsampling:} Mathematically formalized, zoom-aware method computing bucket sizes as functions of wavelength range and zoom level, using overlapping boundaries to prevent peak splitting, and inserting NIST-matched peaks explicitly.

\item \textbf{NIST-based validation:} Automated cross-validation of detected spectral peaks against curated NIST Atomic Spectra Database subsets, with confidence scoring based on wavelength agreement and line strength~\cite{Reader2011}.

\item \textbf{Backend-centric architecture:} Three-tier system (PostgreSQL + Python API + React frontend) where all scientifically meaningful logic resides in auditable, versioned backend modules with MD5 checksums and algorithm version tracking.

\item \textbf{Empirical validation:} Quantitative experiments on actual Chandrayaan-3 L1 files demonstrating cleaning effects, unsaturated zoom ranges, and trade-offs between bucket counts, peak retention, and latency.
\end{enumerate}

\subsection{Paper Organization}

Section~\ref{sec:related} reviews related work. Section~\ref{sec:data} describes the Chandrayaan-3 dataset and PDS4 context. Section~\ref{sec:architecture} presents system architecture. Section~\ref{sec:processing} details data processing and cleaning. Section~\ref{sec:math} provides mathematical formulation of adaptive downsampling. Section~\ref{sec:api} describes API design. Section~\ref{sec:results} presents experimental validation. Section~\ref{sec:discussion} discusses implications and limitations. Section~\ref{sec:conclusion} concludes.

\section{Related Work}
\label{sec:related}

\subsection{Planetary LIBS Instruments and Mission Heritage}

ChemCam on MSL provides remote compositional information using the first planetary LIBS instrument~\cite{Wiens2012,Maurice2012}. SuperCam on Mars 2020 extends this with co-aligned LIBS, Raman, VISIR, and acoustic measurements~\cite{Wiens2020,Bernardi2021}. MarSCoDe on Tianwen-1 combines LIBS and SWIR spectroscopy covering 240--2400~nm~\cite{Xu2021,MarSCoDe2024}. These missions demonstrate LIBS maturity and motivate robust data infrastructures handling diverse instruments and calibration schemes.

\subsection{Planetary Data Systems and Archival Standards}

PDS4 organizes data into bundles, collections, and products using XML labels describing structure, units, and provenance~\cite{Hughes2018}. For LIBS, typical products pair XML with CSV or fixed-width tables containing wavelength, intensity, and flags. While self-describing, these archives lack out-of-the-box APIs or visualization tools~\cite{GDALPDS4}.

\subsection{Downsampling and Feature-Preserving Visualization}

General-purpose time-series downsampling methods (mean/median aggregation, Douglas-Peucker, LTTB) are widely used in web visualization~\cite{amCharts} but prioritize perceived shape over strict peak retention. For spectroscopy, narrow emission lines are primary information carriers. The downsampling problem for interactive, peak-preserving visualization of high-resolution planetary spectra has not been comprehensively addressed~\cite{Cremers2006}.

\section{Chandrayaan-3 LIBS Dataset and PDS4 Context}
\label{sec:data}

\subsection{Mission Overview and Landing Site}

Chandrayaan-3 launched July 14, 2023 and achieved soft landing near the lunar south pole on August 23, 2023~\cite{ISROCh3}. The Pragyan rover carried LIBS for elemental composition and APXS for complementary X-ray spectroscopy. The landing site is of high interest for potential volatile deposits, ancient crustal materials, and future resource utilization.

\subsection{LIBS Level-1 Data Products}

Chandrayaan-3 LIBS data are released as calibrated Level-1 (L1) products conforming to PDS4, consisting of:

\begin{itemize}[nosep]
\item XML label describing observation time, mission/instrument context, and table column definitions;
\item Tabular file (CSV/fixed-width) storing wavelength (\si{nm}), intensity (counts/calibrated units), measurement IDs, timestamps, and flags (laser status, plasma/background).
\end{itemize}

Representative L1 files contain $\sim 2{,}000$ wavelength samples per row spanning 164--878~nm, with multiple measurements per observation (background + plasma shots).

\subsection{PDS4 Semantics and Database Mapping}

PDS4 LIBS products define column semantics, units, and processing levels, but not relational query abstractions. LunarAtlas maps these into a relational schema with tables for mission, instrument, observation, file version, measurement, and spectral samples, preserving provenance (PDS4 logical identifiers, MD5 hashes) while enabling SQL-level queries.

\section{LunarAtlas System Architecture}
\label{sec:architecture}

\subsection{Design Principles}

LunarAtlas follows three core principles: (1)~\textbf{Single source of scientific truth}—all domain logic resides in backend; frontend is rendering only. (2)~\textbf{Reproducibility by design}—every algorithm and parameter is versioned; API responses include versions and commit hashes; each file has MD5 checksum. (3)~\textbf{Separation of concerns}—data persistence, scientific logic, and presentation in distinct tiers.

\subsection{Three-Tier Architecture and Integrity Tracking}

The architecture comprises:

\textbf{Data layer:} PostgreSQL stores mission metadata, measurement tables, spectral samples, NIST reference lines, and MD5 digests.

\textbf{Logic layer:} Python backend (FastAPI/Flask) implements PDS4 parsing, ingestion, cleaning, bucket computation, min-max aggregation, peak detection, NIST matching, MD5 computation, and JSON serialization~\cite{FastAPIIntro}.

\textbf{Presentation layer:} React + amCharts frontend manages viewports and rendering, never performing scientific calculations or client-side downsampling~\cite{amCharts}.

For each ingested file, metadata recorded in \texttt{file\_version} table: original filename/size, MD5 checksum, ingestion timestamp, and algorithm version, enabling traceability to exact L1 product and pipeline version.

\section{Data Processing and Cleaning Pipeline}
\label{sec:processing}

\subsection{Overview of the Python Pipeline}

The pipeline consists of two main scripts:

\textbf{\texttt{batch\_process\_libs.py}:} reads PDS4 L1 CSV files, reshapes wide tables into long-form per-wavelength rows, performs background subtraction and quality filtering, computes MD5 checksums, writes cleaned spectra.

\textbf{\texttt{plot\_libs\_spectra.py}:} loads raw and cleaned spectra, generates diagnostic plots verifying cleaning parameters, exports figures.

\subsection{From Wide L1 Tables to Long-Form Spectra}

Chandrayaan-3 L1 LIBS products store intensities in wide format: initial columns contain metadata (\texttt{Time}, \texttt{Measurement Count}, \texttt{Operation Mode}, \texttt{Measurement Type}, \texttt{Force Reset Status}, \texttt{Laser Fire Status}), followed by wavelength-value columns (e.g., \texttt{164.35}, \texttt{164.74}, \ldots, \texttt{878.26}). The pipeline identifies numeric column names as wavelengths and reshapes into long-form:
\begin{equation}
\texttt{Measurement\_ID},\ \texttt{Time\_UTC},\ \texttt{Measurement\_Count},\ \texttt{Wavelength\_nm},\ \texttt{Raw\_Intensity},\ \ldots
\end{equation}

\subsection{Background Subtraction and Plasma Selection}

For each observation sequence containing plasma and background shots, the pipeline: (1)~identifies valid plasma measurements (\texttt{Is\_Valid\_Plasma = True}, \texttt{Is\_Background = False}); (2)~pairs each plasma spectrum with matched background spectrum; (3)~computes cleaned intensities via
\begin{equation}
I_{\text{cleaned}}(\lambda) = \max\bigl(0,\, I_{\text{plasma}}(\lambda) - I_{\text{background}}(\lambda)\bigr);
\end{equation}
(4)~marks samples with validity flags, retaining instrument metadata.

\subsection{Effect of Cleaning on Chandrayaan-3 Spectrum}

We demonstrate the pipeline effect using one plasma spectrum from \texttt{ch3\_lib\_002\_20230825T104221\_00\_l1.csv}. Wavelength channels span 164.35--878.26~nm.

\begin{figure}[H]
\centering
\includegraphics[width=\textwidth]{libs_raw_vs_cleaned_full.png}
\caption{\textbf{Full-range Chandrayaan-3 LIBS spectrum.} Raw L1 counts (blue) and cleaned intensity (orange) for single plasma measurement spanning 164.35--878.26~nm. Cleaning removes continuum background while preserving emission features.}
\label{fig:raw_clean_full}
\end{figure}

\begin{figure}[H]
\centering
\includegraphics[width=\textwidth]{libs_raw_vs_cleaned_zoom.png}
\caption{\textbf{Zoomed spectrum (380--420~nm).} Comparison of raw L1 counts (blue) with cleaned intensity (orange). Background subtraction removes continuum while preserving emission-line peaks, enhancing signal-to-background ratio.}
\label{fig:raw_clean_zoom}
\end{figure}

\subsection{File Integrity and MD5 Checksums}

At ingestion, the pipeline computes MD5 checksums for every source L1 and exported cleaned file. Recorded metadata: original filename/size, MD5 checksum (128-bit hex), ingestion timestamp, algorithm version and commit hash. This enables reproducibility, traceability, and integrity verification.

\section{Mathematical Formulation of Adaptive Min-Max Downsampling}
\label{sec:math}

\subsection{Bucket Size as a Function of Zoom}

Given viewing window $S = \{(\lambda_i, I_i): i=1,\ldots,N_{\text{raw}}\}$ with $\lambda_{\min} \le \lambda_i \le \lambda_{\max}$, define wavelength span $\Delta\lambda = \lambda_{\max} - \lambda_{\min}$. For zoom level $z\in\mathbb{Z}_{\ge 0}$, unconstrained bucket size is
\begin{equation}
b_{\text{size}}(z) = \frac{\Delta\lambda}{\text{BASE\_BUCKETS}\times 2^z}.
\end{equation}

With $\text{BASE\_BUCKETS}\approx 1000$, this gives roughly one bucket per pixel at $z=0$. To avoid buckets smaller than instrument resolution, minimum bucket size $\text{MIN\_BUCKET\_SIZE}$ (e.g., \SI{0.01}{nm}) is enforced:
\begin{equation}
b_{\text{size,final}} = \max\bigl(b_{\text{size}}(z),\, \text{MIN\_BUCKET\_SIZE}\bigr).
\end{equation}

\subsection{Zoom Saturation Threshold and Valid Zoom Range}

Unsaturated downsampling regime:
\begin{equation}
b_{\text{size}}(z) \ge \text{MIN\_BUCKET\_SIZE} \quad\Longleftrightarrow\quad 2^z \le \frac{\Delta\lambda}{\text{BASE\_BUCKETS}\times \text{MIN\_BUCKET\_SIZE}}.
\end{equation}

Maximum unsaturated zoom level:
\begin{equation}
z_{\max}(\Delta\lambda) = \left\lfloor \log_2\left(\frac{\Delta\lambda}{\text{BASE\_BUCKETS}\times \text{MIN\_BUCKET\_SIZE}}\right) \right\rfloor.
\end{equation}

For $z > z_{\max}$, bucket size clamps to $\text{MIN\_BUCKET\_SIZE}$ and implementation switches to raw-data mode.

\subsection{Worked Example: 300--400~nm Window at $z=6$}

For window $[300, 400]~\text{nm}$ ($\Delta\lambda = 100~\text{nm}$) with $\text{BASE\_BUCKETS} = 1000$, $\text{MIN\_BUCKET\_SIZE} = 0.01~\text{nm}$:
\begin{equation}
b_{\text{size}}(6) = \frac{100}{1000\times 2^6} = \frac{1}{640} \approx 1.56\times 10^{-3}~\text{nm}.
\end{equation}

Applying floor constraint:
\begin{equation}
b_{\text{size,final}} = \max(1.56\times 10^{-3}, 0.01) = 0.01~\text{nm},\quad N_{\text{buckets}} = \left\lceil \frac{100}{0.01} \right\rceil = 10{,}000.
\end{equation}

At $z=6$, algorithm saturates: each bucket is \SI{0.01}{nm} wide; further zoom increases cannot reduce bucket size.

\subsection{Representative Bucket Sizes and $z_{\max}$}

For Chandrayaan-3 spectrum (164.35--878.26~nm, $\Delta\lambda = 713.91$~nm) with $\text{BASE\_BUCKETS} = 1000$, $\text{MIN\_BUCKET\_SIZE} = 0.01$~nm:

\begin{table}[H]
\centering
\caption{\textbf{Zoom saturation analysis.} Bucket sizes and maximum unsaturated zoom levels for representative windows. Beyond $z_{\max}$, system switches to raw data mode.}
\label{tab:zmax}
\small\sffamily
\begin{tabular}{@{}lccc@{}}
\toprule
\textbf{Window} & \boldmath$\Delta\lambda$ \textbf{(nm)} & \boldmath$z_{\max}$ & \boldmath$b_{\text{size}}(z_{\max})$ \textbf{(nm)} \\
\midrule
Full Chandrayaan-3 band & 713.91 & 6 & 0.0112 \\
100~nm window & 100.00 & 3 & 0.0125 \\
10~nm window & 10.00 & 0 & 0.0100 \\
\bottomrule
\end{tabular}
\end{table}

\section{API Design and Visualization Workflow}
\label{sec:api}

\subsection{Spectral Endpoints}

Primary endpoint: \texttt{GET /api/v1/spectral} with query parameters \texttt{observation\_id}, \texttt{min\_wavelength}, \texttt{max\_wavelength}, \texttt{zoom\_level}, \texttt{include\_nist}. Responses include metadata, algorithm version, raw point counts, bucket parameters, and bucket objects with optional NIST annotations.

Additional endpoints provide mission/instrument metadata, measurement catalogs, MD5 digests, and NIST reference lines. All endpoints return machine-readable JSON with explicit units and schema documentation.

\subsection{Progressive Loading and Zoom Strategy}

Frontend interaction: (1)~request overview spectrum (full range, low zoom); (2)~on zoom events, translate viewport to wavelength ranges and zoom levels; (3)~request updated buckets respecting $z_{\max}$ values; (4)~display min-max buckets with optional NIST overlays.

At maximum effective zoom ($z > z_{\max}(\Delta\lambda)$), backend switches to raw mode, streaming all samples for full-resolution inspection.

\subsection{Frontend Visualization}

React + amCharts frontend renders buckets as filled ranges between $I_{\min}$ and $I_{\max}$ versus wavelength, with optional NIST-matched peak markers~\cite{amCharts}. Tooltips display wavelength/intensity ranges, point counts, identified elements with confidence scores. Frontend never performs downsampling or peak detection—pure consumer of backend-computed data.

\section{Results and Experimental Validation}
\label{sec:results}

\subsection{Baseline Suppression and Noise Behaviour}

Analyzing 560--580~nm window for same Chandrayaan-3 measurement: raw L1 spectrum has mean intensity $\sim 1005$~counts, RMS fluctuation $\sim 47$~counts; cleaned spectrum has mean $\sim 141$~counts, RMS $\sim 184$~counts.

This implies: (1)~baseline suppression factor $\sim 7.1$; (2)~high-frequency variance increase factor $\sim 3.9$.

The cleaning strongly reduces low-frequency background while enhancing local contrast—consistent with subtracting smooth continuum and leaving sharper residual structure. For peak detection, baseline reduction is beneficial, increasing peak prominence despite larger local variance.

\subsection{Zoom Range and Peak Preservation}

Using zoom-saturation criterion (Section~\ref{sec:math}), LunarAtlas configures downsampling: for each window, backend computes $z_{\max}(\Delta\lambda)$; for $z \le z_{\max}$, min-max downsampling is active; for $z > z_{\max}$, bucket sizes clamp at $\text{MIN\_BUCKET\_SIZE}$ and API returns raw samples.

Table~\ref{tab:zmax} shows full Chandrayaan-3 band allows unsaturated downsampling to $z_{\max}=6$; 100~nm window saturates at $z_{\max}=3$; 10~nm window already saturated at $z=0$.

\subsection{Performance Metrics}

On modest server configuration: end-to-end response times (query + aggregation + serialization) remain below $\sim \SI{500}{ms}$ for typical L1 spectra across zoom levels; payload sizes 10--400~kB under compression. Composite indices on \texttt{(observation\_id, wavelength)} provide substantial speedups. Min-max buckets reduce rendered points by one to two orders of magnitude while preserving peak amplitudes and positions.

\section{Discussion}
\label{sec:discussion}

LunarAtlas demonstrates that backend-centric architecture with mathematically explicit, peak-preserving downsampling bridges the gap between archival PDS4 data and modern interactive spectral analysis expectations. Centralizing domain logic ensures visualizations are reproducible scientific views, not client-side approximations.

The adaptive min-max approach addresses generic downsampling shortcomings by prioritizing scientifically meaningful extrema and reference lines. Combined with NIST validation and confidence scoring, LunarAtlas democratizes LIBS interpretation while preserving expert-level metadata. The pattern generalizes to other spectroscopic missions (Mars LIBS, Raman, XRF, imaging spectroscopy) with appropriate instrument-specific calibration and reference databases.

\textbf{Limitations and Future Work.} Current limitations: (1)~manual quality assessment of cleaned spectra; (2)~absence of automated continuum fitting and multi-peak deconvolution; (3)~limited NIST coverage for minor/trace elements; (4)~lack of real-time ingestion for ongoing missions. Future work will address these through machine learning-based quality scoring, integration of advanced spectral fitting libraries (\texttt{lmfit}, \texttt{pyspeckit}), expanded reference databases including molecular bands and isotopic lines, and streaming ingestion pipelines for near-real-time mission support.

\section{Conclusion}
\label{sec:conclusion}

LunarAtlas illustrates how publicly available Chandrayaan-3 LIBS data can be transformed into reproducible, interactive, scientifically rigorous spectral infrastructure. Key elements: PDS4-aware ingestion, physically motivated cleaning, mathematically formalized adaptive min-max downsampling tailored to spectroscopy, and backend-centric API exposing peak-preserving, NIST-validated spectra to thin visualization clients.

The system achieves high peak retention across zoom levels, wavelength accuracies within instrumental resolution, and sub-second response times. By treating spectral data, processing algorithms, and derived visualizations as tightly coupled, versioned artefacts, LunarAtlas provides both working implementation for Chandrayaan-3 and generalizable blueprint for future planetary spectroscopy data systems.

As humanity expands lunar, Martian, and beyond exploration, open, reproducible, accessible scientific data systems become increasingly critical. LunarAtlas represents a step toward democratizing planetary science by making mission data not merely available, but truly usable by researchers, educators, and enthusiasts worldwide.

\section*{Acknowledgments}

We acknowledge ISRO for public release of Chandrayaan-3 LIBS data through PDS4-compliant archives, and NIST for maintaining the Atomic Spectra Database underpinning spectral validation~\cite{Reader2011}. We thank Aditya Bhardwaj (Scientist, ISRO National Remote Sensing Centre, and TEDxMAIT speaker) for public talks on India's space programme and real-time satellite data systems, which helped shape LunarAtlas's broader vision. We acknowledge the open-source scientific Python ecosystem (NumPy, SciPy, Pandas, Astropy) and planetary science community whose LIBS instrument and calibration work informs this study.

\section*{Author Contributions}

\textbf{L.~Anand:} Conceptualization, system architecture, algorithm development, database design, API implementation, data analysis, writing—original draft and editing.

\textbf{D.~Saeed:} Data processing pipeline, PDS4 parsing, NIST validation, visualization design, performance benchmarking, writing—review and editing.

Both authors reviewed and approved the final manuscript.

\section*{Data Availability}

Processed Chandrayaan-3 LIBS data products and LunarAtlas source code will be made publicly available upon manuscript acceptance. Raw Chandrayaan-3 LIBS data are available from ISRO's PDS4 archives.

\section*{Competing Interests}

The authors declare no competing interests.

\begin{thebibliography}{99}

\bibitem{Cremers2006}
D.~A. Cremers and L.~J. Radziemski, \textit{Handbook of Laser-Induced Breakdown Spectroscopy}, 2nd ed., John Wiley \& Sons, Chichester, UK, 2013.

\bibitem{Wiens2012}
R.~C. Wiens \textit{et~al.}, ``The ChemCam Instrument Suite on the Mars Science Laboratory (MSL) Rover: Body Unit and Combined System Tests,'' \textit{Space Sci. Rev.}, vol.~170, no.~1--4, pp.~167--227, 2012.

\bibitem{Maurice2012}
S.~Maurice \textit{et~al.}, ``The ChemCam Instrument Suite on the Mars Science Laboratory (MSL) Rover: Science Objectives and Mast Unit Description,'' \textit{Space Sci. Rev.}, vol.~170, no.~1--4, pp.~95--166, 2012.

\bibitem{Mangold2016}
N.~Mangold \textit{et~al.}, ``ChemCam activities and discoveries during the nominal mission of the Mars Science Laboratory in Gale crater, Mars,'' \textit{J. Anal. At. Spectrom.}, vol.~31, no.~4, pp.~863--889, 2016.

\bibitem{Wiens2020}
R.~C. Wiens \textit{et~al.}, ``The SuperCam Instrument Suite on the NASA Mars 2020 Rover: Body Unit and Combined System Tests,'' \textit{Space Sci. Rev.}, vol.~216, article~89, 2020.

\bibitem{Bernardi2021}
P.~Bernardi \textit{et~al.}, ``Optical design and performance of the SuperCam instrument for the Mars 2020 mission,'' in \textit{Proc. SPIE 11852, International Conference on Space Optics — ICSO 2020}, 2021.

\bibitem{Xu2021}
W.~Xu \textit{et~al.}, ``The MarSCoDe Instrument Suite on the Mars Rover of China's Tianwen-1 Mission,'' \textit{Space Sci. Rev.}, vol.~217, article~64, 2021.

\bibitem{MarSCoDe2024}
Y.~Li \textit{et~al.}, ``Classification of Martian rocks and soils at Tianwen-1 landing site based on laser-induced breakdown spectroscopy data of MarSCoDe,'' \textit{Spectrochim. Acta Part B}, vol.~211, article~106836, 2024.

\bibitem{ISROCh3}
Indian Space Research Organisation, ``Chandrayaan-3: Details,'' \url{https://www.isro.gov.in/Chandrayaan3.html}, accessed April 2026.

\bibitem{Hughes2018}
S.~J. Hughes, A.~C. Raugh, and T.~L. Crichton, ``Planetary Data System Standards Reference,'' NASA Planetary Data System, version~1.14.0, 2018.

\bibitem{GDALPDS4}
GDAL/OGR contributors, ``PDS4 -- NASA Planetary Data System (Version 4),'' GDAL documentation, 2024.

\bibitem{Reader2011}
J.~Reader, K.~Olsen, and G.~Nave, ``Atomic Spectroscopic Databases at NIST,'' in \textit{AIP Conf. Proc.}, vol.~1344, pp.~17--26, 2011.

\bibitem{amCharts}
amCharts, ``amCharts: JavaScript Charts \& Maps,'' \url{https://www.amcharts.com/}, accessed April 2026.

\bibitem{FastAPIIntro}
S.~Ramirez, ``FastAPI,'' \url{https://fastapi.tiangolo.com/}, accessed April 2026.

\end{thebibliography}

\end{document}